\chapter{Oscillatory Holes and Phase-Lock Networks}
\label{ch:oscillatory_holes}

\papernote{Source paper: \emph{On the Consequences of Categorical Completion on Cognitive Processes} --- \texttt{docs/categorical-completion/} (oscillatory hole and phase-lock sections)}

\begin{chapterabstract}
As oxygen molecules traverse the \Hplus electric field, they create oscillatory holes: incomplete circuits requiring specific phase patterns to close. Each hole admits $\sim 10^6$ distinct weak-force completions (van der Waals angles, dipole orientations, vibrational phases) that produce the same spatial outcome but different oscillatory signatures. Selecting one configuration from millions \emph{is} thought generation. Neurones possess the unique capacity to generate the required oscillatory patterns; circuit closure \emph{is} conscious experience. This chapter proves the five central corollaries: BMDs are oscillatory holes; thinking requires no more effort than smelling; each conscious moment samples $10^{6 \times 10^{35}}$ completions; holes must be filled (cascade termination causes cell death); and qualia are selection signatures from weak-force arrangements.
\end{chapterabstract}

\input{../docs/categorical-completion/sections/oscillatory-holes.tex}
\input{../docs/categorical-completion/sections/phase-lock-networks.tex}
